\chapter{Plane curves and signed curvature}
\label{chap:curve}

% archon:covers Gluck/Curve.lean

This chapter sets up the elementary differential geometry of plane curves used
throughout the project. We model a plane curve as a map $\gamma : \mathbb{R} \to
\mathbb{C}$, identifying the Euclidean plane $\mathbb{R}^2$ with the complex
plane $\mathbb{C}$; this matches the complex notation $e^{i\theta}$ used in the
source for the reconstruction integral.

\begin{definition}\leanok
[Regular curve]
  \label{def:is_regular}
  \lean{Gluck.IsRegular}
  A map $\gamma : \mathbb{R} \to \mathbb{C}$ is a \emph{regular curve} if it is
  twice continuously differentiable ($C^2$, encoded \texttt{ContDiff ℝ 2}) and
  its velocity never vanishes: $\gamma'(t) \neq 0$ for all $t$.
\end{definition}

% SOURCE: [DeTurck-Gluck], §4, p. 9, lines 163-174 (read from references/deturck-gluck-fourvertex.txt)
% SOURCE QUOTE: "There is a unique map α: [0, 2π] → R2 which begins at the
% origin and has unit tangent vector (cos θ, sin θ) and curvature κ(θ) at the
% point α(θ) . In fact, if s is arc length along this curve, then the equations
% dα/ds = (cos θ, sin θ)  and  dθ/ds = κ(θ)  quickly lead to the explicit formula
% α(θ) = ∫0^θ (cos θ, sin θ) dθ / κ(θ) ."
\begin{definition}\leanok
[Signed curvature]
  \label{def:signed_curvature}
  \lean{Gluck.signedCurvature}
  \uses{def:is_regular}
  \textit{Source: DeTurck--Gluck, §4 (defining relations $d\alpha/ds = (\cos\theta,\sin\theta)$, $d\theta/ds = \kappa$).}
  The \emph{signed curvature} of $\gamma : \mathbb{R} \to \mathbb{C}$ at $t$ is
  \[
    \kappa_\gamma(t) \;=\;
      \frac{\operatorname{Re}\gamma'(t)\,\operatorname{Im}\gamma''(t)
            - \operatorname{Im}\gamma'(t)\,\operatorname{Re}\gamma''(t)}
           {\lvert \gamma'(t)\rvert^{3}} .
  \]
  This is the standard parametrization-invariant formula; for a curve
  parametrized by its angle of inclination $\theta$ (so $\gamma'(\theta) =
  e^{i\theta}/\kappa(\theta)$) it reduces to $\kappa_\gamma(\theta) =
  \kappa(\theta)$, recovering the source's defining relation $d\theta/ds =
  \kappa$.

  \emph{This formula is a $C^2$ notion}: it reads off $\gamma''$, so it is only
  meaningful when $\gamma \in C^2$ (\cref{def:is_regular} bakes in
  \texttt{ContDiff ℝ 2}). The realized curve of a \emph{merely continuous}
  curvature function is only $C^1$, so the final theorem expresses curvature
  realization through the intrinsic notion \cref{def:realizes_curvature} below,
  not through this formula.
\end{definition}

% The intrinsic, C¹-applicable curvature-realization predicate. Project-bespoke
% (no external source beyond the defining relation dθ/ds=κ already quoted for
% def:signed_curvature); introduced to make the final theorem sound for merely
% continuous κ, where the deriv²-formula returns junk.
\begin{definition}

\leanok
[Intrinsic curvature realization]
  \label{def:realizes_curvature}
  \lean{Gluck.RealizesCurvature}
  \uses{def:is_curvature_function}
  A curve $\gamma : \mathbb{R} \to \mathbb{C}$ \emph{realizes the curvature
  function} $\kappa : \mathbb{R} \to \mathbb{R}$ if $\gamma$ is $C^1$, regular
  ($\gamma'(t) \neq 0$ for all $t$), and there is a differentiable
  \emph{tangent-angle} function $\varphi : \mathbb{R} \to \mathbb{R}$ such that,
  for all $t$,
  \[
    \gamma'(t) \;=\; \lVert \gamma'(t)\rVert \, e^{i\varphi(t)},
    \qquad
    \varphi'(t) \;=\; \kappa(t)\,\lVert \gamma'(t)\rVert .
  \]
  The first equation says $\varphi(t)$ is the angle of inclination of the unit
  tangent; the second is the source's defining relation $d\theta/ds = \kappa$
  rewritten via the chain rule with arc length $s$, $ds/dt = \lVert\gamma'(t)
  \rVert$: $\varphi'(t) = (d\varphi/ds)(ds/dt) = \kappa(t)\lVert\gamma'(t)\rVert$.

  \emph{Why this is the right notion.} It needs only $\gamma \in C^1$ (no
  $\gamma''$): the angle $\varphi$ — not $\gamma$ — is required to be
  differentiable, and for the reconstruction curve $\alpha_\rho$ (whose velocity
  is $e^{i\theta}\rho(\theta)$) the tangent angle is exactly the parameter,
  $\varphi(\theta) = \theta$, which is smooth even when $\alpha_\rho$ is only
  $C^1$. Thus $\alpha_\rho$ realizes a merely-continuous $\kappa$. When $\kappa >
  0$, $\varphi' = \kappa\lVert\gamma'\rVert > 0$, so the tangent angle turns
  strictly monotonically: this is the intrinsic form of \emph{convexity}
  (rotation index $1$ over a period), replacing the $C^2$ predicate
  \cref{def:is_convex_curve} in the final theorem.
\end{definition}

\begin{lemma}

\leanok
[Intrinsic realization agrees with the formula on $C^2$ curves]
  \label{lem:realizes_curvature_signedCurvature}
  \lean{Gluck.signedCurvature_of_realizesCurvature}
  \uses{def:realizes_curvature, def:signed_curvature, def:is_regular}
  If $\gamma$ is $C^2$ and regular and realizes $\kappa$
  (\cref{def:realizes_curvature}), then its signed curvature equals $\kappa$:
  $\kappa_\gamma(t) = \kappa(t)$ for all $t$. (Bridge between the intrinsic
  notion and the deriv²-formula; lets the smooth lemma
  \cref{lem:signed_curvature_reconstruct} and the intrinsic realization be used
  interchangeably on $C^2$ data.)
\end{lemma}
\begin{proof}
  \uses{def:realizes_curvature, def:signed_curvature}
  \leanok
  Write $v(t) = \lVert\gamma'(t)\rVert$, so $\gamma'(t) = v(t)e^{i\varphi(t)}$
  with $\varphi$ the tangent angle. Since $\gamma \in C^2$ and $\gamma'$ is
  nonvanishing, $v$ is $C^1$ and $\gamma''(t) = v'(t)e^{i\varphi(t)} + v(t)
  i\varphi'(t)e^{i\varphi(t)}$. Decompose into real/imaginary parts:
  $\Re\gamma' = v\cos\varphi$, $\Im\gamma' = v\sin\varphi$, and
  $\Re\gamma'' = v'\cos\varphi - v\varphi'\sin\varphi$,
  $\Im\gamma'' = v'\sin\varphi + v\varphi'\cos\varphi$. The signed-curvature
  numerator is
  \[
    \Re\gamma'\,\Im\gamma'' - \Im\gamma'\,\Re\gamma''
      = v^2\varphi'(\cos^2\varphi + \sin^2\varphi) = v^2\varphi',
  \]
  (the $v'$ terms cancel), while $\lVert\gamma'\rVert^3 = v^3$. Hence
  $\kappa_\gamma = v^2\varphi'/v^3 = \varphi'/v = \kappa$ by the realization
  relation $\varphi' = \kappa v$.
\end{proof}

\begin{definition}\leanok
[Closed curve]
  \label{def:is_closed_curve}
  \lean{Gluck.IsClosedCurve}
  A curve $\gamma : \mathbb{R} \to \mathbb{C}$ is \emph{closed} if it is
  $2\pi$-periodic: $\gamma(t + 2\pi) = \gamma(t)$ for all $t$.
\end{definition}

\begin{definition}\leanok
[Simple closed curve]
  \label{def:is_simple_closed}
  \lean{Gluck.IsSimpleClosed}
  \uses{def:is_closed_curve}
  A curve $\gamma$ is a \emph{simple closed curve} if it is closed and injective
  on a fundamental period, i.e. $\gamma$ restricted to $[0, 2\pi)$ is injective.
  Such a $\gamma$ is the embedding of the circle $S^1 = \mathbb{R}/2\pi\mathbb{Z}$
  that the converse theorem produces.
\end{definition}

\begin{definition}\leanok
[Convex curve]
  \label{def:is_convex_curve}
  \lean{Gluck.IsConvexCurve}
  \uses{def:signed_curvature}
  A regular curve $\gamma$ is \emph{convex} if its signed curvature is strictly
  positive everywhere: $\kappa_\gamma(t) > 0$ for all $t$. For a simple closed
  curve this is equivalent to convexity of the enclosed region.
\end{definition}

\begin{lemma}\leanok
[Positive curvature forces convexity]
  \label{lem:pos_curvature_convex}
  \lean{Gluck.isConvexCurve_of_signedCurvature_pos}
  \uses{def:is_convex_curve, def:signed_curvature}
  If $\gamma$ is regular and $\kappa_\gamma(t) > 0$ for all $t$, then $\gamma$ is
  convex.
\end{lemma}
\begin{proof}

\leanok
  Immediate from the definition of \cref{def:is_convex_curve}: convexity is
  precisely the everywhere-positivity of the signed curvature.
\end{proof}
